\documentclass[12pt]{article}
\usepackage{../../../lecture_notes}
\usepackage{../../../math}
\usepackage{../../../uark_colors}

\hypersetup{
  colorlinks = true,
  allcolors = ozark_mountains,
  breaklinks = true,
  bookmarksopen = true
}

\begin{document}
\begin{center}
  {\Huge\bf Midterm 1 - Fall 2024}
  
  \smallskip
  {\large\it ECON 4753 — University of Arkansas}
\end{center}

\vspace{5mm}
\begin{enumerate}
  \item This question concerns our general discussion of modelling outcomes. 
  Say we have an outcome variable $Y$ that we want to analyze and a set of variables $X$ we want to use in our model. 
  In general, we write the `true model' as $Y_i = f(X_i) + e_i$.
  With a sample, we will fit some model which we will call $\hat{f}$. 

  Two models were fit using the same training sample.
  Model A is a simple linear regression, while Model B allows a more complicated relationship between $X$ and $Y$.

  \begin{center}
    \begin{tabular}{lc}
      Model & In-sample MSPE \\
      \hline
      A & 420 \\
      B & 310 \\
    \end{tabular}
  \end{center}

  \begin{enumerate}
    \item Write the formula for in-sample mean-squared prediction error.
    
    \item A student argues, ``Model B has a lower in-sample MSPE, so it will predict outcomes more accurately for new observations.''
    Do you agree with this conclusion? Explain.
  \end{enumerate}


  \bigskip\bigskip
  \item This question is based on our review of statistics.
  A company advertises that its wireless earbuds last an average of 7 hours on a single charge.
  You independently test a random sample of $n = 64$ pairs of this model under the same conditions.
  Let $b_i$ denote battery life in hours.
  Your sample has a mean of $\bar{b} = 6.6$ hours and a variance of $\var{b} = 4$ hours squared.
  The sample distribution of the sample mean is therefore $\bar{b} \sim \mathcal{N}( \mu, 1 / 16 )$, where $\mu$ is the true average battery life.

  \medskip
  \begin{enumerate}[label=(\alph*)]
    \item Form a 95\% confidence interval for the population mean battery life.
    Interpret this interval in words.

    \item You suspect that the average battery life is less than the advertised 7 hours.
    Test the null hypothesis that $H_0: \mu = 7$ at the 5\% significance level.
  \end{enumerate}


  %% Spotify questions
  \newpage
  \item The following questions will focus on a dataset from Spotify. 
  A random sample of songs were drawn from 6 genres: `EDM', `latin', `r\&b', `pop', `rap', and `rock'.
  Spotify produced a few variables we will use for the analysis: (1) \texttt{track\_popularity} measures track populartiy on a scale from 0 to 100 and (2) \texttt{danceability} measures how suitable a track is for dancing on a scale from 0 to 1.

  I ran a regression of track popularity on the song's danceability measure. 

  \begin{codeblock}[{}]
OLS estimation, Dep. Var.: track_popularity
Observations: 32,833
             Estimate Std. Error t value  Pr(>|t|)    
(Intercept)   35.1757   0.632777 55.5894 < 2.2e-16 ***
danceability  11.1497   0.945354 11.7942 < 2.2e-16 ***
---
Signif. codes:  0 '***' 0.001 '**' 0.01 '*' 0.05 '.' 0.1 ' ' 1
  \end{codeblock}

  \begin{enumerate}
    \item Interpret the coefficient on a song's \texttt{danceability} measure
    
    \item Form a 95\% confidence interval for the coefficient. 
  \end{enumerate}
  

  \newpage

  \vspace{5mm}
  \item Continuing with the Spotify dataset, we want to see if different genre's have different popularities. We regress the track's popularity on indicators for each genre (excluding one)
  \begin{codeblock}[{}]
OLS estimation, Dep. Var.: track_popularity
Observations: 32,833
               Estimate Std. Error   t value  Pr(>|t|)    
(Intercept)   47.744870   0.339019 140.83227 < 2.2e-16 ***
genre::edm   -12.911344   0.451279 -28.61054 < 2.2e-16 ***
genre::latin  -0.718294   0.490232  -1.46521   0.14287    
genre::r&b    -6.521339   0.488259 -13.35632 < 2.2e-16 ***
genre::rap    -4.529416   0.457639  -9.89736 < 2.2e-16 ***
genre::rock   -6.016532   0.489295 -12.29632 < 2.2e-16 ***
---
Signif. codes:  0 '***' 0.001 '**' 0.01 '*' 0.05 '.' 0.1 ' ' 1
  \end{codeblock}

  \begin{enumerate}
    \item What is the omitted group in this regression?

    \item Interpret the coefficient on rock in words.
    
    \item Which genre is more popular, r\&b or rap?

    \item How would you modify this regression to test whether this difference is statistically significant?
  \end{enumerate}
\end{enumerate}



\end{document}
